% index_extraction_list.tex
\documentclass[12pt]{article}
\usepackage[utf8]{inputenc}
\usepackage[paper=a4paper,margin=2.5cm]{geometry}
\begin{document}

\section*{אינדקס extraction\_list.tex}
\textbf{מטרה}:
להגדיר את שדות החילוץ הסטנדרטיים ואת פורמט הפלט.

\subsection*{שדות מרכזיים}
\begin{itemize}
  \item \textbf{Source ID} מזהה מקור ייחודי.
  \item \textbf{Title} כותרת המקור.
  \item \textbf{Authors} שמות המחברים.
  \item \textbf{Year} שנת פרסום.
  \item \textbf{Abstract Summary} תקציר 1-3 משפטים.
  \item \textbf{Methodology} שיטת המחקר.
  \item \textbf{Sample Size} גודל מדגם.
  \item \textbf{Key Findings} ממצאים מרכזיים.
  \item \textbf{Limitations} מגבלות.
  \item \textbf{Data Availability} זמינות נתונים וקישורים.
  \item \textbf{Quality Score} דירוג איכות.
  \item \textbf{Extracted Tables} טבלאות נחוצות.
  \item \textbf{Notes} הערות חופשיות.
\end{itemize}

\subsection*{פורמט פלט מומלץ}
המלצה על JSON עם שדות תואמים לשם אינטגרציה עם כלי ניתוח.
\end{document}
